\documentclass[a4paper,12pt]{article}
\usepackage[T2A]{fontenc}
\usepackage[utf8]{inputenc}
\usepackage[russian]{babel}
\usepackage{amsmath,amssymb,mathtools}
\usepackage{PTSans}
\renewcommand{\familydefault}{\sfdefault}
\usepackage{geometry}
\geometry{left=23mm,right=23mm,top=20mm,bottom=22mm}
\usepackage{microtype}
\usepackage{array,tabularx,booktabs}
\usepackage{xcolor}
\usepackage{tikz}
\usepackage{hyperref}
\usepackage{titlesec}
\usepackage{needspace}

\definecolor{accent}{HTML}{006E70}
\definecolor{ink}{HTML}{252D32}
\definecolor{muted}{HTML}{606B70}
\definecolor{critical}{HTML}{A12C31}
\hypersetup{colorlinks=true,linkcolor=accent,
  pdfauthor={Лёвин Артём Александрович},
  pdftitle={Ревью домашней работы — Анна Баннова, 8 ноября 2026}}
\titleformat{\section}{\Large\bfseries\color{ink}}{}{0pt}{}
\titlespacing*{\section}{0pt}{0pt}{1.2ex}
\titleformat{\subsection}{\normalfont\bfseries\normalsize\color{ink}}{}{0pt}{}
\titlespacing*{\subsection}{0pt}{1.5ex plus .3ex}{.35ex}
\setlength{\parindent}{0pt}
\setlength{\parskip}{4pt}
\renewcommand{\arraystretch}{1.30}
\setlength{\emergencystretch}{2em}
\pagestyle{plain}

\newcommand{\errorlabel}{\textcolor{critical}{\textbf{Ошибка:}}}
\newcommand{\keylabel}{\textcolor{accent}{\textbf{Ключевая идея:}}}
\newcommand{\resultlabel}{\textcolor{accent}{\textbf{Итог:}}}
\newcommand{\footerline}{%
  \par\vfill\begin{center}\small\color{muted}
  Лёвин Артём Александрович эксклюзивно для Анны Банновой
  \end{center}}
\begin{document}
\pagenumbering{roman}
\begin{titlepage}
\centering
\vspace*{21mm}
{\large\color{muted} МАТЕМАТИКА\par}
\vspace{20mm}
{\Huge\bfseries\color{ink} Ревью домашней\\[5pt] работы\par}
\vspace{12mm}
{\large Задачи на движение: таблицы,\\ скорости и время старта\par}
\vspace{19mm}
{\Large Анна Баннова\par}
\vspace{8mm}
{\large 8 ноября 2026 г.\par}
\vfill
{\color{muted}\normalsize Лёвин Артём Александрович\\[4pt]
эксклюзивно для Анны Банновой\par}
\vspace{13mm}
\begin{tikzpicture}
\draw[accent!75,line width=.9pt] (-4,0) .. controls (-2.1,.32) and (-.7,-.23) .. (1.1,0) .. controls (2.4,.19) and (3.1,.14) .. (4,0);
\end{tikzpicture}
\end{titlepage}
\pagenumbering{arabic}

\section*{Сводная таблица}
\vspace{3mm}
\small
\begin{tabularx}{\linewidth}{@{}>{\raggedright\arraybackslash}p{13mm}>{\raggedright\arraybackslash}p{28mm}X>{\raggedright\arraybackslash}p{20mm}>{\raggedright\arraybackslash}p{21mm}@{}}
\toprule
\textbf{№} & \textbf{Тема} & \textbf{Суть ошибки} & \textbf{Ваш ответ} & \textbf{Верный ответ}\\
\midrule
\hyperref[task:8]{8} & Движение с разницей во времени & Скорость грузовика записана больше скорости автобуса; последующие обозначения противоречат таблице. & $20$ км/ч & $20$ км/ч\\
\addlinespace[1mm]
\hyperref[task:10]{10} & Поворот и задержка старта & Разность времён взята в обратном порядке: учтён неверный знак задержки. & $3{,}3$ км & $2$ км\\
\bottomrule
\end{tabularx}
\normalsize
\vspace{7mm}
{\small\color{muted} В таблицу включены только задания с ошибками в решении или обосновании.}

\newpage
\thispagestyle{empty}
\vspace*{\fill}
\begin{center}
{\Huge\bfseries\color{ink} Разбор ошибок}
\end{center}
\vspace*{\fill}
\newpage

\phantomsection\label{task:8}
\section*{Задание № 8}
\subsection*{Текст задания}
Автобус и грузовик проехали по 24 км. Грузовик двигался на 5 км/ч медленнее автобуса, но выехал на 24 минуты раньше, поэтому прибыли они одновременно. Найдите скорость автобуса.

\subsection*{Ваше решение}
Вы перевели 24 минуты в часы: $24/60=0{,}4$ ч. Затем обозначили скорость автобуса через $y$, а в строке грузовика указали скорость $y+5$. В дальнейших вычислениях появляется также выражение $y-5$. После преобразований получено квадратное уравнение
\[
 y^2-5y-300=0,
\]
и выбран корень $y=20$.

\subsection*{Ваш ответ}
$20$ км/ч. \textbf{Число правильное, но приведённые обозначения и переходы не образуют корректного решения.}

\subsection*{Ошибка}
\errorlabel\ В исходной таблице нарушен смысл слов «грузовик двигался на 5 км/ч медленнее автобуса». Если автобус едет со скоростью $y$, то грузовик должен ехать со скоростью $y-5$, а в таблице указано $y+5$.

\subsection*{Где именно возникает ошибка}
Последние надёжно правильные данные Вашего решения:
\[
 S=24\text{ км},\qquad 24\text{ мин}=\frac{24}{60}=0{,}4\text{ ч}.
\]
Первый неправильный шаг --- назначение скорости грузовика:
\[
 \underbrace{v_{\text{авт}}=y}_{\text{правильно}},\qquad
 \underbrace{v_{\text{груз}}=y+5}_{\text{неверно}}.
\]
В последующих формулах одновременно встречаются $y+5$ и $y-5$. При одном и том же обозначении $y$ это разные модели движения. В частности, записанная разность $24/(y+5)-24/y$ при $y>0$ отрицательна и не может равняться положительным $0{,}4$ часа.

\subsection*{Почему это неверно}
Если грузовик медленнее автобуса на 5 км/ч, его скорость \emph{меньше} скорости автобуса на 5. Для одной и той же дистанции медленный участник тратит больше времени, поскольку время равно расстоянию, делённому на скорость:
\[
 t=\frac{S}{v}.
\]
Грузовик выехал раньше и прибыл одновременно с автобусом, поэтому находился в пути на $0{,}4$ часа дольше. Требуется вычитать \emph{время автобуса из времени грузовика}.

\subsection*{Ключевая идея}
\keylabel\ Сначала определите, кто быстрее и кто дольше ехал. Только после этого составьте таблицу и запишите разность времён с правильным знаком:
\[
 t_{\text{груз}}-t_{\text{авт}}=0{,}4.
\]

\subsection*{Исправление}
Сохраняем выбранное Вами обозначение $y$ для скорости автобуса и правильный перевод минут в часы. Исправляем строку грузовика:
\[
 v_{\text{авт}}=y,\qquad v_{\text{груз}}=y-5,\qquad y>5.
\]
Отсюда сразу получается уравнение
\[
 \frac{24}{y-5}-\frac{24}{y}=0{,}4.
\]

\subsection*{Полное правильное решение}
Пусть $y$ км/ч --- скорость автобуса. Тогда скорость грузовика $(y-5)$ км/ч. По смыслу задачи $y>5$, иначе скорость грузовика была бы неположительной.

\begin{center}
\small
\begin{tabular}{@{}lccc@{}}
\toprule
 & Скорость, км/ч & Расстояние, км & Время, ч\\
\midrule
Автобус & $y$ & $24$ & $\dfrac{24}{y}$\\[2mm]
Грузовик & $y-5$ & $24$ & $\dfrac{24}{y-5}$\\
\bottomrule
\end{tabular}
\end{center}

Оба прибыли одновременно, но грузовик стартовал на $0{,}4$ часа раньше. Значит, разница между временами в пути равна $0{,}4$ часа:
\[
\begin{aligned}
\frac{24}{y-5}-\frac{24}{y}&=0{,}4,\\[2mm]
\frac{24y-24(y-5)}{y(y-5)}&=0{,}4,\\[2mm]
\frac{120}{y(y-5)}&=0{,}4,\\
120&=0{,}4y(y-5),\\
300&=y(y-5),\\
y^2-5y-300&=0.
\end{aligned}
\]
Дискриминант:
\[
D=(-5)^2-4\cdot1\cdot(-300)=25+1200=1225=35^2.
\]
Корни:
\[
y_{1,2}=\frac{5\pm35}{2},\qquad y_1=20,\quad y_2=-15.
\]
Корень $-15$ не удовлетворяет условию $y>5$. Следовательно, скорость автобуса --- $20$ км/ч.

\subsection*{Проверка}
Скорость грузовика равна $20-5=15$ км/ч. Его время в пути:
\[
\frac{24}{15}=1{,}6\text{ ч}=96\text{ мин}.
\]
Автобус ехал
\[
\frac{24}{20}=1{,}2\text{ ч}=72\text{ мин}.
\]
Разница $96-72=24$ минуты, как в условии.

\subsection*{Итог}
\resultlabel\ \textbf{Скорость автобуса $20$ км/ч.} Ответ совпал с Вашим, но исходную таблицу и уравнение времени необходимо согласовать с условиями «медленнее» и «раньше».

\subsection*{Задача на закрепление}
Автобус и грузовик проехали по 18 км. Грузовик двигался на 6 км/ч медленнее автобуса, но отправился на 30 минут раньше. Оба прибыли одновременно. Найдите скорость автобуса.

\footerline
\newpage

\phantomsection\label{task:10}
\section*{Задание № 10}
\subsection*{Текст задания}
От старта до опушки 8 км. Медленный турист начал путь со скоростью 4 км/ч. Через 15 минут вслед за ним стартовал быстрый турист со скоростью 8 км/ч, дошёл до опушки и повернул обратно. Туристы встретились, пока медленный ещё шёл к опушке. На каком расстоянии от опушки произошла встреча? Ответ дайте в километрах.

\subsection*{Ваше решение}
Вы обозначили через $x$ расстояние от опушки до точки встречи. В таблице верно записали пройденные пути: медленный турист --- $8-x$, быстрый турист --- $8+x$. Времена Вы выразили правильно:
\[
 t_{\text{медл}}=\frac{8-x}{4},\qquad
 t_{\text{быстр}}=\frac{8+x}{8}.
\]
После этого составлено уравнение
\[
 \frac{8+x}{8}-\frac{8-x}{4}=\frac{15}{60},
\]
из которого получено $x=10/3\approx3{,}3$ км.

\subsection*{Ваш ответ}
Приблизительно $3{,}3$ км.

\subsection*{Ошибка}
\errorlabel\ В уравнении перепутан порядок вычитания времён. Быстрый турист отправился на 15 минут \emph{позже}, поэтому до встречи он ехал на 15 минут \emph{меньше} медленного туриста. Нужно вычесть время быстрого туриста из времени медленного.

\subsection*{Где именно возникает ошибка}
Последний правильный шаг --- формулы для расстояний и времён:
\[
\frac{8-x}{4}\quad\text{и}\quad\frac{8+x}{8}.
\]
Первое неверное равенство:
\[
\underbrace{\frac{8+x}{8}-\frac{8-x}{4}=\frac14}_{\text{разность записана с неверным знаком}}.
\]
В этом равенстве утверждается, что быстрый турист находился в пути на четверть часа \emph{дольше}, хотя по условию он стартовал позже. Для одновременной встречи верно противоположное соотношение.

\subsection*{Почему этот шаг неверен}
Пусть медленный турист отправился в момент начала отсчёта времени. Быстрый отправился через $15$ минут, то есть через $\tfrac14$ часа. Если встреча произошла спустя $T$ часов после старта медленного, медленный двигался $T$ часов, а быстрый --- $T-\tfrac14$ часа. Поэтому
\[
 T-\left(T-\frac14\right)=\frac14.
\]
Общий принцип: \emph{время того, кто вышел раньше, больше на величину задержки}, если участники встретились в один и тот же момент.

\subsection*{Ключевая идея}
\keylabel\ При движении с поворотом путь быстрого туриста до встречи складывается из пути до опушки и пути назад. А задержка старта отражается в \emph{разности времён}, причём знак определяется тем, кто начал раньше:
\[
 \frac{8-x}{4}-\frac{8+x}{8}=\frac14.
\]

\subsection*{Исправление}
Сохраняем Ваши правильные выражения для путей $8-x$ и $8+x$. Меняем только порядок вычитания времён:
\[
\begin{aligned}
\frac{8-x}{4}-\frac{8+x}{8}&=\frac14,\\[1mm]
\frac{2(8-x)-(8+x)}8&=\frac14,\\
2(8-x)-(8+x)&=2,\\
16-2x-8-x&=2,\\
8-3x&=2,\\
3x&=6,\\
x&=2.
\end{aligned}
\]

\subsection*{Полное правильное решение}
Пусть $x$ км --- расстояние от опушки до места встречи. Поскольку медленный турист ещё не дошёл до опушки, $0<x<8$.

Медленный турист прошёл от старта до места встречи $8-x$ км. При скорости 4 км/ч время движения равно
\[
 t_{\text{медл}}=\frac{8-x}{4}\text{ ч}.
\]
Быстрый турист сначала прошёл 8 км до опушки, затем повернул назад и прошёл ещё $x$ км. Его полный путь до встречи --- $8+x$ км, а время движения равно
\[
 t_{\text{быстр}}=\frac{8+x}{8}\text{ ч}.
\]
Быстрый отправился через 15 минут после медленного, но встретились они одновременно. Значит, медленный находился в пути дольше на $15/60=1/4$ часа:
\[
\frac{8-x}{4}-\frac{8+x}{8}=\frac14.
\]
Приводим дроби к общему знаменателю и раскрываем скобки:
\[
\begin{aligned}
2(8-x)-(8+x)&=2,\\
16-2x-8-x&=2,\\
8-3x&=2,\\
x&=2.
\end{aligned}
\]
Полученное значение удовлетворяет ограничению $0<x<8$.

\subsection*{Проверка}
Точка встречи находится в $8-2=6$ км от старта. Медленный турист шёл
\[
 \frac64=1{,}5\text{ ч}=90\text{ мин}.
\]
Быстрый прошёл $8+2=10$ км, затратив
\[
 \frac{10}{8}=1{,}25\text{ ч}=75\text{ мин}.
\]
Разница $90-75=15$ минут. К моменту встречи быстрый уже дошёл до опушки и вернулся на 2 км, а медленный ещё находился в 2 км от неё. Все условия выполняются.

Проверим также Ваше значение $x=10/3$. Тогда медленный находился бы в пути $7/6$ часа, а быстрый --- $17/12$ часа: быстрый ехал бы \emph{дольше} на четверть часа, что противоречит более позднему старту.

\subsection*{Итог}
\resultlabel\ \textbf{Встреча произошла в $2$ км от опушки.} Расстояния Вы выразили правильно; исправлению подлежит знак разности времён.

\subsection*{Задача на закрепление}
От старта до озера 12 км. Первый турист отправился к озеру со скоростью 4 км/ч. Через 30 минут вслед за ним отправился второй турист со скоростью 6 км/ч. Второй дошёл до озера, повернул обратно и встретился с первым, когда тот ещё шёл к озеру. На каком расстоянии от озера произошла встреча?

\footerline
\end{document}
